
Domain mixing ratios can substantially affect LLM downstream performance. DoReMi demonstrated that optimized mixing yields 2-3\% improvements over heuristic baselines on benchmarks including MMLU. SlimPajama showed that careful curation of domain ratios is essential for efficient pretraining. These findings establish that training data composition matters as much as data quantity.

Current optimization methods carry computational overhead. DoReMi requires training both a reference model and a proxy model before determining optimal weights---a process that can consume 10-20\% of total training compute. For practitioners under compute constraints, this overhead represents a barrier to optimizing data mixtures.

Existing approaches fundamentally require training to evaluate domain utility. Perplexity-based selection needs model inference; proxy training needs gradient updates. No method currently predicts domain utility in a training-free manner.

This gap motivates our investigation: can embedding geometry---the similarity between domain samples and downstream task exemplars---provide a signal for domain utility without training? Prior work in domain adaptation suggests that distributional alignment correlates with transfer performance. If embedding space proximity to task-relevant content predicts training utility, domain mixtures could potentially be optimized using only frozen embedder inference.

We propose Embedding-guided Domain Mixing Prediction (EDMP), a training-free approach that scores domains by cosine similarity between domain sample embeddings and task exemplar embeddings.

Our contributions are:

\begin{enumerate}
\item \textbf{Pipeline Validation:} We demonstrate that EDMP scores can be reliably computed for multi-domain corpora, producing statistically distinguishable domain rankings (ANOVA F=1242.59, p<0.001) with perfect reproducibility across seeds.
\end{enumerate}

\begin{enumerate}
\item \textbf{Data Quality Dependency:} We identify that synthetic data with shared vocabulary produces insufficient cross-domain variance (standard deviation 0.0069 versus target 0.05), establishing that real domain data is necessary to test the predictive power hypothesis.
\end{enumerate}

\begin{enumerate}
\item \textbf{Methodological Foundation:} We provide a reusable pipeline and experimental framework for embedding-based domain scoring.
\end{enumerate}

The core hypothesis---that EDMP scores predict downstream training utility---was not tested in this work due to data limitations.

